\documentclass[10pt,a4paper]{amsart}
\usepackage[margin=19mm]{geometry}
\usepackage{amsmath,amssymb,mathtools}
\usepackage[scr=rsfs,cal=euler]{mathalfa}
\usepackage{xfrac}
\usepackage[unicode,hidelinks,bookmarks=false]{hyperref}
\newcommand{\ie}{i.e.\ }
\newcommand{\cf}{cf.\ }
\newcommand{\resp}{respectively\ }
\newcommand{\Subsection}{\subsection}
\providecommand{\nobreakdash}{\nobreak\hbox{-}}
\setlength{\emergencystretch}{2em}
\multlinegap0pt
\usepackage{amssymb}
\newtheorem{definition}{Definition}
\newtheorem{theorem}{Theorem}
\title{On planar sections of the dodecahedron}
\author{Andreas Thom}
\date{Source: arXiv:2511.12773v4 (CC BY 4.0)}
\begin{document}
\maketitle

\noindent\emph{Source and licence.} On planar sections of the dodecahedron. Version arXiv:2511.12773v4 (CC BY 4.0), distributed by arXiv under the Creative Commons Attribution 4.0 International licence, http://creativecommons.org/licenses/by/4.0/, as recorded in the arXiv licence metadata for this exact version and shown on the abstract page https://arxiv.org/abs/2511.12773v4. Full licence notice: \texttt{licenses/arxiv-2511.12773-CC-BY-4.0.txt}.\par\smallskip

\noindent\textit{Editorial scope.} Counted unit: the article's theorem thm:convex-bodies (source0.tex lines 260--266). The article's complete original proof of it is delivered with it (source0.tex lines 268--393, one \texttt{proof} environment of 126 lines). No mathematics is written, repaired or re-derived: the statement and the proof are the article's own lines, in the article's own notation, and no retained line is substituted or re-flowed.  Retained uncounted as the source-local context the proof needs: lines 174--179; lines 192--194.  Nothing was omitted from any retained span, so the package carries no omission record.  No retained line contains a citation, so the wrapper carries no bibliography and no bibliography entry.  No retained line contains a figure environment, an image-inclusion command or drawing code, so no image is delivered and the package declares no asset dependency.  Editorial formatting only, in addition: an AMS \texttt{amsart} preamble that restates the article's own declarations and macros used by the retained lines, namely the environments \texttt{definition}, \texttt{theorem} on separate counters, together with the standard packages those lines need.  The retained environments print in this document as Definition 1, Theorem 1 in order of appearance, whereas the article prints the same items with the numbering of its own full document.  The complete native source of the article as distributed in the arXiv e-print for this exact version is bundled unchanged as \texttt{sources/189228/source0.tex}.
\begin{definition}\label{def:planar-statistics}
The \emph{planar statistic} of $X\subset V$ is the
multiset $\mathrm{PS}(X)$
of planar congruence classes $[\Pi]_X$ as $\Pi$ ranges over all vertex-planes
in $\mathcal{H}$, counted with multiplicity. 
\end{definition}

\begin{theorem}\label{thm:nonuniq}
There exist two 7-element subsets $S,T\subset V$ that are not congruent in $\mathbb{R}^3$, yet $\mathrm{PS}(S)=\mathrm{PS}(T)$. One such pair is $S=\{0,1,2,3,4,11,17\}$ and $T=\{0,1,3,4,5,11,17\}$.
\end{theorem}

\begin{theorem}\label{thm:convex-bodies}
There exist three-dimensional polytopes $K,L \subseteq \mathbb R^3$, such that
\begin{enumerate}
\item $K$ and $L$ are not congruent, but
\item $(\Phi_K)_*(\lambda) = (\Phi_L)_*(\lambda)$.
\end{enumerate}
\end{theorem}

\begin{proof}
Let $D\subset\mathbb{R}^3$ be a regular dodecahedron with vertex set $V$, and
let $S,T\subset V$ be two subsets as in Theorem~\ref{thm:nonuniq}: they
are not congruent in $\mathbb{R}^3$, but their planar statistics agree, i.e.\ 
$\mathrm{PS}(S) = \mathrm{PS}(T)$
as multisets of planar congruence classes in the sense of
Definition~\ref{def:planar-statistics}.


Let us now explain how we construct polytopes out of $S$ and $T$. For each $X\subset V$ define a convex polytope $K_X$ by truncating all vertices
$v\in X$ by small congruent triangular caps. We will show that $K_S$ and $K_T$ are not
congruent in $\mathbb{R}^3$, but the distributions of isometry classes of their
planar sections with respect to the invariant measure on affine planes coincide.

Fix once and for all a small parameter $\varepsilon>0$.
For each vertex $v\in V$ let $P_v$ be the unique plane that intersects the three
edges incident to $v$ in points of distance $\varepsilon$ from $v$, and denote
by $H_v^-$ the closed half-space bounded by $P_v$ that contains the barycenter
of $D$. For any subset $X\subset V$ we set
\[
K_X := D\cap\bigcap_{v\in X} H_v^-.
\]
By construction $K_X$ is obtained from $D$ by cutting off, at each $v\in X$, a
small congruent triangular pyramid. In particular, if $g\in \mathrm{Isom}(\mathbb{R}^3)$ is an
isometry with $g(K_S)=K_T$, then $g$ must map the set of truncated vertices of
$K_S$ onto that of $K_T$, and restricts to an isometry $g\colon D\to D$ with
$g(S)=T$. Since $S$ and $T$ are not congruent as subsets of $V$ by
Theorem~\ref{thm:nonuniq}, it follows that $K_S$ and $K_T$ are not
congruent.

We now compare the distributions of planar sections. For
$H\in\mathcal{A}(3,2)$ set
\[
m(H) := \bigl|\{v\in V : H\cap B(v,\varepsilon)\neq\emptyset\}\bigr|,
\]
the number of vertex neighbourhoods met by $H$. For each integer $m\ge 0$ define
\[
A_m := \{H\in\mathcal{A}(3,2) : m(H)=m\},
\]
so that $\mathcal{A}(3,2) = \sqcup_{m\ge 0} A_m.$
Let $p_m := \lambda(A_m)$ and $\lambda_m := p_m^{-1}\,\lambda\!\restriction_{A_m}$
for all $m$ with $p_m>0$. Then, we have
\[
\lambda = \sum_{m\ge 0} p_m\,\lambda_m,
\]
a canonical convex decomposition of $\lambda$, and for each $X\subset V$ we
obtain a corresponding convex decomposition
\[
\mu_X = \sum_{m\ge 0} p_m\,\mu_{X,m},
\qquad
\mu_{X,m} := (\Phi_{K_X})_*(\lambda_m).
\]
To prove $\mu_S=\mu_T$ it therefore suffices to show that $\mu_{S,m} = \mu_{T,m}$ for all $m\ge 0$.
\medskip

\emph{The cases $m=0$ and $1$.}
If $H\in A_0$, then $H$ misses all balls $B(v,\varepsilon)$ and in particular all caps.
The section $K_X\cap H$ then coincides with $D\cap H$; it does not depend on $X$. In particular, we can conclude $\mu_{S,0} = \mu_{T,0}.$
Similarly, if $H\in A_1$, then $H$ meets exactly one ball $B(v,\varepsilon)$, say at $v$. The local
shape of the section $K_X\cap H$ in a neighbourhood of $B(v,\varepsilon)$ is completely
determined by the local geometry at $v$ and thus only depends on whether $v\in X$ or not. Thus, since $D$
is vertex-transitive and $|S|=|T|$, we obtain $\mu_{S,1} = \mu_{T,1}$.

\emph{The case $m\ge 3$.}
Suppose $H\in A_m$ with $m\ge 3$. By choice of $\varepsilon>0$, this implies that
$H$ lies in a small neighbourhood of a unique vertex plane
$\Pi\subset\mathbb{R}^3$ with
\[
\Pi\cap V = \{v_1,\dots,v_m\},\qquad m\ge 3,
\]
and that the vertices $\{v_1,\dots,v_m\}$ are precisely those whose balls
$B(v_i,\varepsilon)$ meet $H$. The local configuration of the section $K_X\cap H$ is then,
up to isometry, completely determined by the isometry type of the inclusion $\Pi\cap X \subset V$
and the position of $H$ in a sufficiently small neighbourhood of $\Pi$.

The isometry group $D$ acts on the set of such local configurations. For
a fixed $X \subset V$ the orbits of this action are parametrized precisely by
the planar congruence classes $[\Pi]_X$ appearing in the planar statistic
$\mathrm{PS}(X)$: two vertex planes $\Pi,\Pi'$ give rise to the same local
configuration type if and only if there is an isometry $g$ such that $g(V)=V$ and
$g(\Pi\cap X)=\Pi'\cap X$.

Fix a $G$-orbit of local configurations $\mathcal{O}$, i.e.\ an orbit of a planar inclusion $Z \subset V$ with $|Z|\ge 3$. Consider the set of planes $H$ whose local configuration belongs to $\mathcal{O}$. By
invariance of $\lambda$ and congruence of all caps and faces, the restriction of
$\lambda$ to this set has the same total mass and the same normalized image
measure under $\Phi_X$, independent of $X$; it depends only on $\mathcal{O}$.
Thus for each such orbit $\mathcal{O}$ there is a probability measure
$\nu_{\mathcal{O}}$ on $\mathcal{P}$ describing the contribution of one
occurrence of $\mathcal{O}$, and the component $\mu_{X,m}$ can be written as a
finite convex combination
\[
\mu_{X,m} = \sum_{\mathcal{O}} c_{\mathcal{O}}(X)\,\nu_{\mathcal{O}},
\]
where the sum runs over all orbits $\mathcal{O}$ with $m\ge 3$, and
$c_{\mathcal{O}}(X)$ is a nonnegative coefficient proportional to the number of
vertex planes whose local configuration with respect to $X$ lies in $\mathcal{O}$.
By definition, the multiset of classes $[\Pi]_X$ in $\mathrm{PS}(X)$ records
exactly these multiplicities. Thus the coefficients $c_{\mathcal{O}}(X)$ are
completely determined by $\mathrm{PS}(X)$.

Since $\mathrm{PS}(S) = \mathrm{PS}(T)$ as multisets of planar congruence
classes, it follows that $c_{\mathcal{O}}(S) = c_{\mathcal{O}}(T)$ for every
orbit $\mathcal{O}$, and hence $\mu_{S,m} = \mu_{T,m}$ for all $m \geq 3$.

\emph{The cases $m=2$.} If $H\in A_2$, then $H$ meets exactly two balls $B(v,\varepsilon)$, say at vertices
$v,w\in V$. The local shape of $K_X\cap H$ is then, up to isometry, determined by
the orbit type of the unordered pair $\{v,w\}$ under the isometry group of
$D$ together with the membership of $v,w$ in $X$. Again, there are only finitely many
such orbit types of unordered pairs in $V$, and for each such type all
corresponding pairs have the same local geometry. Each such pair belongs to various vertex planes with three or more vertices, so its contribution is again
encoded in the planar statistic $\mathrm{PS}(X)$; in particular, the counts of
pairs of each orbit type contained in $X$ are determined by $\mathrm{PS}(X)$.
Since $\mathrm{PS}(S)=\mathrm{PS}(T)$, the contributions for $m=2$ agree: $\mu_{S,2} = \mu_{T,2}$.

\medskip

Combining all cases, we have shown that $\mu_{S,m} = \mu_{T,m}$ for every $m\ge 0$. Since
\[
\mu_S = \sum_{m\ge 0} p_m\,\mu_{S,m}
\quad\text{and}\quad
\mu_T = \sum_{m\ge 0} p_m\,\mu_{T,m}
\]
with the same weights $p_m$, we conclude that $\mu_S=\mu_T$. This shows that the
distributions of isometry classes of planar sections of $K_S$ and $K_T$ with
respect to $\lambda$ coincide, as claimed.
\end{proof}

\end{document}
